\section{Introduction}

Low-latency execution of deep neural networks (DNN) plays a critical role in autonomous driving \cite{cordts2016cityscapes}, augmented reality \cite{alhaija2017augmented}, language translation\cite{devlin2018bert}, and other applications of AI.
DNNs can be expressed as a directed acyclic computational graph (DAG), in which nodes represent the operators (\textit{e.g.}, convolution, matrix multiplication) and directed edges represent the dependencies between operators.
Existing deep learning frameworks (\textit{e.g.}, Tensorflow~\cite{abadi2016tensorflow}, PyTorch~\cite{paszke2019pytorch}, MXNet~\cite{chen2015mxnet}) map the operators in DNNs to vendor-provided kernel libraries (\textit{e.g.}, cuDNN~\cite{chetlur2014cudnn}, MKL-DNN~\cite{intel2017mkldnn}) to achieve high performance.
However, these kernel libraries require significant engineering effort to manually tune for each hardware platform and operator.
The significant manual effort required to produce efficient operator implementations for each target accelerator limits the development and innovation of new operators~\cite{barham2019machine} and specialized accelerators~\cite{moreau2019hardware}.

Given the importance of DNNs' performance, researchers and industry practitioners have turned to search-based compilation~\cite{chen2018tvm, adams2019learning,zheng2020flextensor,vasilache2018tensor,liu2019optimizing} for automated generation of \emph{tensor programs}, \textit{i.e.}, low-level implementations of tensor operators. 
For an operator or a (sub-)graph of multiple operators, users define the computation in a high-level declarative language (\autoref{sec:background}), and the compiler then searches for programs tailored towards different hardware platforms.

To find performant tensor programs, it is necessary for a search-based approach to explore a large enough search space to cover all the useful tensor program optimizations. 
However, existing approaches fail to capture many effective optimization combinations, because they rely on either predefined manually-written templates (\textit{e.g.}, TVM~\cite{chen2018learning}, FlexTensor~\cite{zheng2020flextensor}) or aggressive pruning by evaluating incomplete programs (\textit{e.g.}, Halide auto-scheduler \cite{adams2019learning}), which prevents them from covering a comprehensive search space (\autoref{sec:background}).
The rules they use to construct the search space are also limited.

In this paper, we explore a novel search strategy for generating high-performance tensor programs.
It can automatically generate a large search space with comprehensive coverage of optimizations and gives every tensor program in the space a chance to be chosen. 
It thus enables to find high-performance programs that existing approaches miss.

Realizing this goal faces multiple challenges.
First, it requires automatically constructing a large search space to cover as many tensor programs as possible for a given computation definition.
Second, we need to search efficiently without comparing incomplete programs in the large search space that can be orders of magnitude larger than what existing templates can cover.
Finally, when optimizing an entire DNN with many subgraphs, we should recognize and prioritize the subgraphs that are critical to the end-to-end performance.

To this end, we design and implement \emph{\Ansor}, a framework for automated tensor program generation.
\Ansor utilizes a hierarchical representation to cover a large search space.
This representation decouples high-level structures and low-level details, enabling flexible enumeration of high-level structures and efficient sampling of low-level details.
The space is constructed automatically for a given computation definition.
\Ansor then samples complete programs from the search space and fine-tunes these programs with evolutionary search and a learned cost model.
To optimize the performance of DNNs with multiple subgraphs, 
\Ansor dynamically prioritizes subgraphs of the DNNs that are more likely to improve the end-to-end performance.


We evaluate \Ansor on both standard deep learning benchmarks and emerging new workloads against manual libraries and state-of-the-art search-based frameworks.
Experiment results show that \Ansor improves the execution performance of DNNs on the Intel CPU, ARM CPU, and NVIDIA GPU by up to  $3.8\times$, $2.6\times$, and $1.7\times$, respectively.
For most computation definitions, the best program found by \Ansor is outside the search space of existing search-based approaches.
The results also show that, compared with existing search-based approaches, \Ansor searches more efficiently, generating higher-performance programs in a shorter time, despite its larger search space.
\Ansor can match the performance of a state-of-the-art framework with an order of magnitude less search time.
Besides, \Ansor enables automatic extension to new operators by only requiring their mathematical definitions without manual templates.

In summary, this paper makes the following contributions:

\begin{itemize}[topsep=-0.1em, itemsep=-0.1em]
    \item A mechanism to generate a large hierarchical search space of tensor programs for a computational graph.
    \item An evolutionary strategy with a learned cost model to fine-tune the performance of tensor programs.
    \item A scheduling algorithm based on gradient descent to prioritize important subgraphs when optimizing the end-to-end performance of DNNs.
    \item An implementation and comprehensive evaluation of the \Ansor system demonstrating that the above techniques outperform state-of-the-art systems on a variety of DNNs and hardware platforms.
\end{itemize}
